\part{Foundations}

\chapter{Why build an LLM inference engine?}
\label{ch:why-build}
An inference engine is a mathematical model wrapped in a file-format contract, a runtime architecture, and a validation discipline. This book starts with GPT-2 because it is compact enough to audit yet exercises the core decoder-only path. \citep{vaswani2017attention,radford2018gpt,radford2019gpt2}



\begin{intuitionbox}{Intuition}
The first useful engine is the one whose behaviour is explainable and reproducible across runtimes.
\end{intuitionbox}

\begin{definitionbox}{Mathematical or contract statement}
For token ids $x_{1:T}$, inference computes logits $z_T=f_\theta(x_{1:T})\in\R^V$. The project target is tolerance-based agreement across PyTorch, C++, and Rust over the same converted artifacts.
\end{definitionbox}

\begin{implementationbox}{Implementation interpretation}
Build the narrow complete slice first: converter, PyTorch reference, golden vectors, C++ engine, Rust engine, shared validation, then benchmark baseline. CUDA is studied early but implemented later.
\end{implementationbox}

\begin{verificationbox}{How to verify this works}
\begin{itemize}
\item Check shapes and dtypes at the boundary.
\item Compare deterministic outputs against PyTorch golden vectors.
\item Record tolerance, seed, model artifact, and backend.
\item Do not benchmark until validation passes.
\end{itemize}
\end{verificationbox}

\begin{warningbox}{Common failure modes}
\begin{itemize}
\item Letting framework defaults choose dtype or layout.
\item Testing only a batch size of one.
\item Depending on upstream checkpoint names in runtime code.
\item Treating benchmark output as validation.
\end{itemize}
\end{warningbox}

\begin{chapterexercises}
\item State the central invariant in one sentence.\\
\textit{Hint.} Use the conceptual guide vocabulary.
\item Construct a toy example with small shapes.\\
\textit{Hint.} Use $B=2$, $T=3$, $H=8$ where possible.
\item Name a validation test that would catch a plausible bug.\\
\textit{Hint.} Prefer generated golden vectors to handwritten long arrays.
\end{chapterexercises}

\chapter{Mathematical notation for tensors and shapes}
\label{ch:tensors-shapes}
Most inference bugs begin as shape bugs. This chapter fixes the notation for tensors, broadcasting, layout, and cache dimensions.


\begin{table}[htbp]\centering\small
\begin{tabularx}{\textwidth}{L{0.20\textwidth}L{0.20\textwidth}Y}
\toprule
Symbol & Toy value & Meaning \\
\midrule
$B$ & 1 or 2 & Batch size. \\
$T$ & 3 or 4 & Current input sequence length. \\
$S$ & accumulated & Visible cache sequence length. \\
$L$ & 2 or 12 & Number of transformer blocks. \\
$P$ & 16 or 1024 & Maximum learned-position table length. \\
$H$ & 8 or 16 & Hidden size. \\
$F$ & 32 or 64 & Feed-forward intermediate width. \\
$H_q$ & 2 or 4 & Query attention heads. \\
$H_{kv}$ & $H_q$ in phase 0 & Key/value heads. \\
$d_h$ & $H/H_q$ & Per-head dimension. \\
$g$ & $H_q/H_{kv}$ & Query heads per key/value head for GQA. \\
$V$ & toy vocab or 50257 & Vocabulary size. \\
\bottomrule
\end{tabularx}
\caption{Shape notation used throughout the book.}
\end{table}


\begin{intuitionbox}{Intuition}
A shape is a precondition, not a comment.
\end{intuitionbox}

\begin{definitionbox}{Mathematical or contract statement}
Token ids are \(I\in\N^{B\times T}\). Hidden states are \(X\in\R^{B\times T\times H}\). Query tensors are \(Q\in\R^{B\times H_q\times T\times d_h}\). For layer \(\ell\), cached keys and values are logically \(K_\ell,V_\ell\in\R^{B\times H_{kv}\times S\times d_h}\). A causal mask \(M\in\R^{T\times S}\) uses finite values for visible keys and \(-\infty\) for hidden keys.
\end{definitionbox}

\begin{implementationbox}{Implementation interpretation}
Store shape, dtype, device, and layout metadata on tensor objects. Validate \(H \bmod H_q=0\), \(H_q \bmod H_{kv}=0\), and cache sequence bounds before execution. Keep logical layout distinct from physical storage.
\end{implementationbox}

\begin{verificationbox}{How to verify this works}
\begin{itemize}
\item Check shapes and dtypes at the boundary.
\item Compare deterministic outputs against PyTorch golden vectors.
\item Record tolerance, seed, model artifact, and backend.
\item Do not benchmark until validation passes.
\end{itemize}
\end{verificationbox}

\begin{warningbox}{Common failure modes}
\begin{itemize}
\item Letting framework defaults choose dtype or layout.
\item Testing only a batch size of one.
\item Depending on upstream checkpoint names in runtime code.
\item Treating benchmark output as validation.
\end{itemize}
\end{warningbox}

\begin{chapterexercises}
\item State the central invariant in one sentence.\\
\textit{Hint.} Use the conceptual guide vocabulary.
\item Construct a toy example with small shapes.\\
\textit{Hint.} Use $B=2$, $T=3$, $H=8$ where possible.
\item Name a validation test that would catch a plausible bug.\\
\textit{Hint.} Prefer generated golden vectors to handwritten long arrays.
\end{chapterexercises}

\chapter{Linear algebra for inference engineers}
\label{ch:linear-algebra}
Affine maps and matrix multiplication power attention projections, MLPs, and the final language-modelling head.



\begin{intuitionbox}{Intuition}
Algebra tells us what to compute; layout and bandwidth tell us what it costs.
\end{intuitionbox}

\begin{definitionbox}{Mathematical or contract statement}
For \(X\in\R^{N\times H}\), \(W\in\R^{O\times H}\), and \(b\in\R^O\), a linear layer computes \(Y=XW^\top+\mathbf{1}_N b^\top\in\R^{N\times O}\). For a hidden tensor \(X\in\R^{B\times T\times H}\), use \(N=BT\) or apply the same affine map independently to each \((b,t)\) row.
\end{definitionbox}

\begin{implementationbox}{Implementation interpretation}
Implement a clear CPU reference first. Use non-square toy matrices so transposition errors cannot hide.
\end{implementationbox}

\begin{verificationbox}{How to verify this works}
\begin{itemize}
\item Check shapes and dtypes at the boundary.
\item Compare deterministic outputs against PyTorch golden vectors.
\item Record tolerance, seed, model artifact, and backend.
\item Do not benchmark until validation passes.
\end{itemize}
\end{verificationbox}

\begin{warningbox}{Common failure modes}
\begin{itemize}
\item Letting framework defaults choose dtype or layout.
\item Testing only a batch size of one.
\item Depending on upstream checkpoint names in runtime code.
\item Treating benchmark output as validation.
\end{itemize}
\end{warningbox}

\begin{chapterexercises}
\item State the central invariant in one sentence.\\
\textit{Hint.} Use the conceptual guide vocabulary.
\item Construct a toy example with small shapes.\\
\textit{Hint.} Use $B=2$, $T=3$, $H=8$ where possible.
\item Name a validation test that would catch a plausible bug.\\
\textit{Hint.} Prefer generated golden vectors to handwritten long arrays.
\end{chapterexercises}

\chapter{Probability, logits, softmax, and sampling}
\label{ch:probability-sampling}
The model produces logits. Samplers turn logits into token ids using greedy, temperature, top-k, or top-p policies.



\begin{intuitionbox}{Intuition}
Logits are scores; softmax is normalisation; sampling is policy.
\end{intuitionbox}

\begin{definitionbox}{Mathematical or contract statement}
Stable softmax uses \(m=\max_i z_i\) and \(p_i=\exp(z_i-m)/\sum_j\exp(z_j-m)\). Greedy sampling returns \(\arg\max_i z_i\). Temperature sampling applies \(\operatorname{softmax}(z/\tau)\) for \(\tau>0\). Top-\(k\) masks all but the \(k\) largest logits; top-\(p\) keeps the smallest descending-probability prefix whose cumulative probability is at least \(p\).
\end{definitionbox}

\begin{implementationbox}{Implementation interpretation}
Validate logits and greedy token ids for model math. Stochastic sampling needs explicit seed, filter order, and categorical draw.
\end{implementationbox}

\begin{verificationbox}{How to verify this works}
\begin{itemize}
\item Check shapes and dtypes at the boundary.
\item Compare deterministic outputs against PyTorch golden vectors.
\item Record tolerance, seed, model artifact, and backend.
\item Do not benchmark until validation passes.
\end{itemize}
\end{verificationbox}

\begin{warningbox}{Common failure modes}
\begin{itemize}
\item Letting framework defaults choose dtype or layout.
\item Testing only a batch size of one.
\item Depending on upstream checkpoint names in runtime code.
\item Treating benchmark output as validation.
\end{itemize}
\end{warningbox}

\begin{chapterexercises}
\item State the central invariant in one sentence.\\
\textit{Hint.} Use the conceptual guide vocabulary.
\item Construct a toy example with small shapes.\\
\textit{Hint.} Use $B=2$, $T=3$, $H=8$ where possible.
\item Name a validation test that would catch a plausible bug.\\
\textit{Hint.} Prefer generated golden vectors to handwritten long arrays.
\end{chapterexercises}

\chapter{Numerical precision and error tolerances}
\label{ch:precision}
\phasezero{} artifacts are \dtype{float32}. Later half precision, bfloat16, TensorFloat-32, or quantisation require explicit new contracts.



\begin{intuitionbox}{Intuition}
Floating-point equality is a test policy, not a wish.
\end{intuitionbox}

\begin{definitionbox}{Mathematical or contract statement}
A tensor comparison passes when \(|a_i-b_i|\leq \text{atol}+\text{rtol}|b_i|\) for all entries after shape and dtype checks pass. Validators should report \(\max_i |a_i-b_i|\), the largest relative error over nonzero references, and the first failing index. Phase-0 hidden states and logits use explicit tolerances, while greedy ids match exactly.
\end{definitionbox}

\begin{implementationbox}{Implementation interpretation}
Carry dtype through loaders, tensors, ops, exporters, and validators. Report max absolute error, max relative error, and first failing index.
\end{implementationbox}

\begin{verificationbox}{How to verify this works}
\begin{itemize}
\item Check shapes and dtypes at the boundary.
\item Compare deterministic outputs against PyTorch golden vectors.
\item Record tolerance, seed, model artifact, and backend.
\item Do not benchmark until validation passes.
\end{itemize}
\end{verificationbox}

\begin{warningbox}{Common failure modes}
\begin{itemize}
\item Letting framework defaults choose dtype or layout.
\item Testing only a batch size of one.
\item Depending on upstream checkpoint names in runtime code.
\item Treating benchmark output as validation.
\end{itemize}
\end{warningbox}

\begin{chapterexercises}
\item State the central invariant in one sentence.\\
\textit{Hint.} Use the conceptual guide vocabulary.
\item Construct a toy example with small shapes.\\
\textit{Hint.} Use $B=2$, $T=3$, $H=8$ where possible.
\item Name a validation test that would catch a plausible bug.\\
\textit{Hint.} Prefer generated golden vectors to handwritten long arrays.
\end{chapterexercises}

\chapter{PyTorch for reference implementations}
\label{ch:pytorch}
PyTorch is the phase-0 correctness oracle and golden-vector exporter, but it must consume the same internal artifact format as C++ and Rust. \citep{pytorchDocs}



\begin{intuitionbox}{Intuition}
The best reference implementation is boring and inspectable.
\end{intuitionbox}

\begin{definitionbox}{Mathematical or contract statement}
The reference exports selected tensors: input ids, position ids, embedding output, block outputs, final hidden, logits, prefill logits, decode logits, and generated ids.
\end{definitionbox}

\begin{implementationbox}{Implementation interpretation}
Construct PyTorch modules from internal tensor names. Do not validate against an upstream model object that bypasses artifact conversion.
\end{implementationbox}

\begin{verificationbox}{How to verify this works}
\begin{itemize}
\item Check shapes and dtypes at the boundary.
\item Compare deterministic outputs against PyTorch golden vectors.
\item Record tolerance, seed, model artifact, and backend.
\item Do not benchmark until validation passes.
\end{itemize}
\end{verificationbox}

\begin{warningbox}{Common failure modes}
\begin{itemize}
\item Letting framework defaults choose dtype or layout.
\item Testing only a batch size of one.
\item Depending on upstream checkpoint names in runtime code.
\item Treating benchmark output as validation.
\end{itemize}
\end{warningbox}

\begin{chapterexercises}
\item State the central invariant in one sentence.\\
\textit{Hint.} Use the conceptual guide vocabulary.
\item Construct a toy example with small shapes.\\
\textit{Hint.} Use $B=2$, $T=3$, $H=8$ where possible.
\item Name a validation test that would catch a plausible bug.\\
\textit{Hint.} Prefer generated golden vectors to handwritten long arrays.
\end{chapterexercises}

\chapter{Rust for safe systems inference}
\label{ch:rust}
Rust expresses ownership of buffers, tensor views, memory maps, and cache state. \citep{rustBook,rust2024edition}



\begin{intuitionbox}{Intuition}
Make invalid memory states hard to represent.
\end{intuitionbox}

\begin{definitionbox}{Mathematical or contract statement}
A Rust tensor view should carry dtype, shape, and a lifetime tied to backing storage. Fallible artifact loading returns typed errors.
\end{definitionbox}

\begin{implementationbox}{Implementation interpretation}
Organise crates around core contracts, ops, model adapters, CLI, and tests. Avoid unsafe kernels until the safe correctness path is proven.
\end{implementationbox}

\begin{verificationbox}{How to verify this works}
\begin{itemize}
\item Check shapes and dtypes at the boundary.
\item Compare deterministic outputs against PyTorch golden vectors.
\item Record tolerance, seed, model artifact, and backend.
\item Do not benchmark until validation passes.
\end{itemize}
\end{verificationbox}

\begin{warningbox}{Common failure modes}
\begin{itemize}
\item Letting framework defaults choose dtype or layout.
\item Testing only a batch size of one.
\item Depending on upstream checkpoint names in runtime code.
\item Treating benchmark output as validation.
\end{itemize}
\end{warningbox}

\begin{chapterexercises}
\item State the central invariant in one sentence.\\
\textit{Hint.} Use the conceptual guide vocabulary.
\item Construct a toy example with small shapes.\\
\textit{Hint.} Use $B=2$, $T=3$, $H=8$ where possible.
\item Name a validation test that would catch a plausible bug.\\
\textit{Hint.} Prefer generated golden vectors to handwritten long arrays.
\end{chapterexercises}

\chapter{C++ for portable high-performance inference}
\label{ch:cpp}
C++20 or later is a natural fit for CPU correctness and later high-performance paths. \citep{cpp2020,cmakeDocs}



\begin{intuitionbox}{Intuition}
C++ should be explicit before it is clever.
\end{intuitionbox}

\begin{definitionbox}{Mathematical or contract statement}
A C++ tensor view is a pointer or span plus shape, dtype, layout, and owner context. RAII owns buffers and mappings.
\end{definitionbox}

\begin{implementationbox}{Implementation interpretation}
Keep scalar CPU kernels simple. Add SIMD, threading, and CUDA interop only after PyTorch parity.
\end{implementationbox}

\begin{verificationbox}{How to verify this works}
\begin{itemize}
\item Check shapes and dtypes at the boundary.
\item Compare deterministic outputs against PyTorch golden vectors.
\item Record tolerance, seed, model artifact, and backend.
\item Do not benchmark until validation passes.
\end{itemize}
\end{verificationbox}

\begin{warningbox}{Common failure modes}
\begin{itemize}
\item Letting framework defaults choose dtype or layout.
\item Testing only a batch size of one.
\item Depending on upstream checkpoint names in runtime code.
\item Treating benchmark output as validation.
\end{itemize}
\end{warningbox}

\begin{chapterexercises}
\item State the central invariant in one sentence.\\
\textit{Hint.} Use the conceptual guide vocabulary.
\item Construct a toy example with small shapes.\\
\textit{Hint.} Use $B=2$, $T=3$, $H=8$ where possible.
\item Name a validation test that would catch a plausible bug.\\
\textit{Hint.} Prefer generated golden vectors to handwritten long arrays.
\end{chapterexercises}

\chapter{CUDA programming model, introduced early}
\label{ch:cuda-intro}
CUDA is not a phase-0 deliverable, but its execution and memory model should shape interface design from the start. \citep{cudaProgrammingGuide,cudaArchive,cuda13ReleaseNotes}

\begin{figure}[htbp]
\centering
\resizebox{0.94\textwidth}{!}{%
\begin{tikzpicture}[node distance=8mm,>=Latex]
\tikzstyle{level}=[rounded rectangle,draw=EngineBlue!80,fill=EngineBlue!6,minimum width=28mm,minimum height=8mm,align=center,font=\small]
\node[level] (grid) {grid};
\node[level,below left=of grid] (block1) {block};
\node[level,below right=of grid] (block2) {block};
\node[level,below=of block1] (warp1) {warp};
\node[level,below=of block2] (warp2) {warp};
\node[level,below=of warp1] (thread1) {threads};
\node[level,below=of warp2] (thread2) {threads};
\draw[->,thick] (grid)--(block1);
\draw[->,thick] (grid)--(block2);
\draw[->,thick] (block1)--(warp1);
\draw[->,thick] (block2)--(warp2);
\draw[->,thick] (warp1)--(thread1);
\draw[->,thick] (warp2)--(thread2);
\node[rounded rectangle,draw=EngineGreen!80,fill=EngineGreen!7,align=left,text width=44mm,right=15mm of block2,font=\small] {registers\\shared memory\\global memory\\constant memory};
\end{tikzpicture}
}
\caption{CUDA execution and memory hierarchy.}
\end{figure}


\begin{intuitionbox}{Intuition}
A GPU is a throughput machine whose performance depends on parallelism and data movement.
\end{intuitionbox}

\begin{definitionbox}{Mathematical or contract statement}
CUDA launches grids of blocks of threads. Warps execute together, and memory hierarchy includes registers, shared memory, global memory, and constant memory.
\end{definitionbox}

\begin{implementationbox}{Implementation interpretation}
Keep \code{Device}, layout, and profiling concepts in the design even while phase 0 is CPU-only.
\end{implementationbox}

\begin{verificationbox}{How to verify this works}
\begin{itemize}
\item Check shapes and dtypes at the boundary.
\item Compare deterministic outputs against PyTorch golden vectors.
\item Record tolerance, seed, model artifact, and backend.
\item Do not benchmark until validation passes.
\end{itemize}
\end{verificationbox}

\begin{warningbox}{Common failure modes}
\begin{itemize}
\item Letting framework defaults choose dtype or layout.
\item Testing only a batch size of one.
\item Depending on upstream checkpoint names in runtime code.
\item Treating benchmark output as validation.
\end{itemize}
\end{warningbox}

\begin{chapterexercises}
\item State the central invariant in one sentence.\\
\textit{Hint.} Use the conceptual guide vocabulary.
\item Construct a toy example with small shapes.\\
\textit{Hint.} Use $B=2$, $T=3$, $H=8$ where possible.
\item Name a validation test that would catch a plausible bug.\\
\textit{Hint.} Prefer generated golden vectors to handwritten long arrays.
\end{chapterexercises}

\chapter{Testing, verification, golden vectors, and reproducibility}
\label{ch:testing}
Validation is part of the engine contract. PyTorch exports vectors; C++ and Rust consume them.



\begin{intuitionbox}{Intuition}
A good vector localises a bug.
\end{intuitionbox}

\begin{definitionbox}{Mathematical or contract statement}
JSON smoke vectors and NPZ tensor dumps carry model family, model name, input ids, dtype, tolerance, and named arrays. Model-math validation is separate from tokenizer validation.
\end{definitionbox}

\begin{implementationbox}{Implementation interpretation}
Generate vectors from PyTorch; do not paste long arrays into LaTeX or source comments. Benchmark suites are separate from correctness vectors.
\end{implementationbox}

\begin{verificationbox}{How to verify this works}
\begin{itemize}
\item Check shapes and dtypes at the boundary.
\item Compare deterministic outputs against PyTorch golden vectors.
\item Record tolerance, seed, model artifact, and backend.
\item Do not benchmark until validation passes.
\end{itemize}
\end{verificationbox}

\begin{warningbox}{Common failure modes}
\begin{itemize}
\item Letting framework defaults choose dtype or layout.
\item Testing only a batch size of one.
\item Depending on upstream checkpoint names in runtime code.
\item Treating benchmark output as validation.
\end{itemize}
\end{warningbox}

\begin{chapterexercises}
\item State the central invariant in one sentence.\\
\textit{Hint.} Use the conceptual guide vocabulary.
\item Construct a toy example with small shapes.\\
\textit{Hint.} Use $B=2$, $T=3$, $H=8$ where possible.
\item Name a validation test that would catch a plausible bug.\\
\textit{Hint.} Prefer generated golden vectors to handwritten long arrays.
\end{chapterexercises}
